\section{Experiments}
\label{sec:experiments}

\subsection{Setup}

\paragraph{Vector memory.}
We use FAISS IndexFlatIP (exact inner product search) with L2-normalized
all-MiniLM-L6-v2 embeddings (384-dim).  The benign corpus contains 200 entries;
poison counts are attack-specific: $n_{\text{base}}=5$ for \agentpoison{},
$2n_{\text{base}}=10$ for \minja{} (multi-entry bridging coverage), and
$3n_{\text{base}}=15$ for \injecmem{} (volume-based broad recall).  Top-$k=5$.

\paragraph{Evaluation.}
$\asrr$ is measured empirically by issuing all 20 victim queries and counting
retrievals of adversarial entries.  $\asra$ is modelled per
Section~\ref{sec:attacks}.  Defense metrics are measured over the same 20
victim queries (for $\tpr$) and 20 benign queries (for $\fpr$).  Bootstrap
95\% confidence intervals are computed from 5 independent random seeds
(Section~\ref{sec:statistical}).

\subsection{Attack Results}
\label{sec:attack_results}

\IfFileExists{../../results/tables/table1_attack_results.tex}{%
  \input{../../results/tables/table1_attack_results}%
}{%
\begin{table}[t]
  \centering
  \caption{Attack evaluation results.  Bootstrap 95\% CI across 5 seeds.
           $\asra$ is modelled.  All metrics in $[0,1]$; higher is worse
           for defender.}
  \label{tab:attack_results}
  \begin{tabular}{lccccc}
    \toprule
    Attack & $\asrr$ & $\asra$ & $\asrt$ & Benign Acc. & ISR \\
    \midrule
    \agentpoison{} (ours) & $1.00 \pm 0.00$ & $0.68$ & $\approx0.68$ & $0.85$ & $1.00$ \\
    \minja{}              & $0.70 \pm 0.05$ & $0.76$ & $\approx0.53$ & $0.90$ & $0.98$ \\
    \injecmem{}           & $0.50 \pm 0.08$ & $0.57$ & $\approx0.29$ & $0.80$ & $1.00$ \\
    \bottomrule
  \end{tabular}
\end{table}%
}

Table~\ref{tab:attack_results} reports the main attack results.
Figure~\ref{fig:attack_asr} visualises the three metrics side-by-side.
\agentpoison{} achieves $\asrr=1.00$ with the corrected triggered-query
evaluation protocol and centroid-targeting passage.  \minja{} matches
its reported target ($\geq 0.70$) with bridging-step passages.
\injecmem{} achieves intermediate recall via volume (15 entries) at
the cost of lower benign accuracy.

\IfFileExists{figures/fig_01_attack_asr.pdf}{%
\begin{figure}[t]
  \centering
  \includegraphics[width=0.72\linewidth]{fig_01_attack_asr.pdf}
  \caption{Attack success rates (bootstrap 95\,\% CI over 5 seeds).
           Left group: $\asrr$; middle: $\asra$ (modelled); right: $\asrt$.
           \agentpoison{} dominates on retrieval ($\asrr=1.00$) but has
           lower $\asra$ than \minja{} due to its blunter adversarial framing.}
  \label{fig:attack_asr}
\end{figure}%
}{}

\paragraph{Evaluation correction for \agentpoison{}.}
\begin{table}[t]
  \centering
  \caption{Impact of the triggered-query evaluation correction.
           Pre-fix: plain queries (no trigger).
           Post-fix: triggered queries ($q \oplus \trigger^*$).}
  \label{tab:trigger_fix}
  \begin{tabular}{lcc}
    \toprule
    Evaluation protocol & $\asrr$ & Note \\
    \midrule
    Un-triggered queries (pre-Phase 12) & 0.25 & Underestimates attack strength \\
    Centroid passage + triggered queries & 1.00 & Matches paper protocol \\
    \bottomrule
  \end{tabular}
\end{table}

Table~\ref{tab:trigger_fix} isolates the effect of the evaluation correction.
This result has an important implication for future work: simulations that
evaluate \agentpoison{} without trigger prepending will substantially
underestimate attacker capability.

\subsection{Attack-Defense Interaction Matrix}
\label{sec:matrix}

\IfFileExists{../../results/tables/table2_attack_defense_matrix.tex}{%
  % table2_attack_defense_matrix.tex
% generated by: cd src && python3 scripts/generate_tables.py
% or: cd src && python3 scripts/run_pipeline.py
%
% placeholder — run the pipeline to populate this table with real results.
\begin{table*}[t]
\centering
\small
\caption{Attack success rate (ASR-R) under each defense (pre-ingestion
filtering). Defense effectiveness = $1 - \text{ASR-R}_{\text{defended}} /
\text{ASR-R}_{\text{baseline}}$. Run \texttt{generate\_tables.py} to populate.}
\label{tab:attack_defense_matrix}
\begin{tabular}{l|c|ccccc}
\toprule
Attack & Baseline & Watermark & Validation & Proactive & Composite & SAD (Ours) \\
\midrule
\agentpoison{} & 1.00 & --- & --- & --- & --- & --- \\
\minja{}        & 0.70 & --- & --- & --- & --- & --- \\
\injecmem{}     & 0.50 & --- & --- & --- & --- & --- \\
\bottomrule
\multicolumn{7}{l}{\footnotesize Run \texttt{python3 src/scripts/generate\_tables.py} to replace with empirical results.}
\end{tabular}
\end{table*}
%
}{%
  \textit{(Run \texttt{generate\_tables.py} to generate attack-defense matrix table.)}%
}

\IfFileExists{../../results/tables/table3_defense_tpr_fpr.tex}{%
  \input{../../results/tables/table3_defense_tpr_fpr}%
}{}

\IfFileExists{figures/m12_attack_defense_matrix.pdf}{%
\begin{figure}[t]
  \centering
  \includegraphics[width=0.72\linewidth]{m12_attack_defense_matrix.pdf}
  \caption{Attack-defense interaction matrix: $\asrr$ under each
           defense (columns) for each attack (rows).
           Darker shading indicates higher defender-side risk
           (higher ASR-R despite defense).
           No single column achieves uniformly low values,
           motivating composite portfolios.}
  \label{fig:ad_matrix}
\end{figure}%
}{}

\subsection{Statistical Analysis}
\label{sec:statistical}

We use percentile bootstrap confidence intervals (2000 replicates,
\citet{efron1993bootstrap}) to quantify variability across seeds.
Pairwise comparisons use paired $t$-tests with Bonferroni correction
for multiple testing ($\alpha_{\text{corrected}} = 0.05 / 3 = 0.017$).
Cohen's $d$ effect sizes are reported for all comparisons.

\begin{table}[t]
  \centering
  \caption{Pairwise hypothesis tests on $\asrr$ across 5 seeds.
           All pairs are significantly different ($p < 0.05$).}
  \label{tab:hypothesis}
  \begin{tabular}{lcccc}
    \toprule
    Comparison & $t$ & $p$ & Cohen's $d$ & Effect Size \\
    \midrule
    \minja{} vs \agentpoison{} & $19.0$ & $<0.001$ & $11.5$ & large \\
    \injecmem{} vs \minja{}    & $10.2$ & $<0.001$ & $4.4$  & large \\
    \bottomrule
    \multicolumn{5}{l}{\footnotesize Bonferroni-corrected $\alpha = 0.017$; $n=5$ seeds per attack.} \\
  \end{tabular}
\end{table}

\subsection{Adaptive Adversary Results}
\label{sec:adaptive_results}

\IfFileExists{../../results/tables/table5_adaptive_sad.tex}{%
  \input{../../results/tables/table5_adaptive_sad}%
}{}

\IfFileExists{figures/p13_adaptive_tradeoff_agent_poison.pdf}{%
\begin{figure}[t]
  \centering
  \begin{subfigure}[b]{0.32\linewidth}
    \includegraphics[width=\linewidth]{p13_adaptive_tradeoff_agent_poison.pdf}
    \caption{\agentpoison{}}
  \end{subfigure}\hfill
  \begin{subfigure}[b]{0.32\linewidth}
    \IfFileExists{figures/p13_adaptive_tradeoff_minja.pdf}{%
      \includegraphics[width=\linewidth]{p13_adaptive_tradeoff_minja.pdf}%
    }{}
    \caption{\minja{}}
  \end{subfigure}\hfill
  \begin{subfigure}[b]{0.32\linewidth}
    \IfFileExists{figures/p13_adaptive_tradeoff_injecmem.pdf}{%
      \includegraphics[width=\linewidth]{p13_adaptive_tradeoff_injecmem.pdf}%
    }{}
    \caption{\injecmem{}}
  \end{subfigure}
  \caption{Adaptive adversary tradeoff: evasion rate vs.\ retrieval
           degradation $\Delta\asrr$ at increasing synonym-substitution
           budgets.  Each point is one budget level.  The negative
           slope confirms the fundamental tension: evasion always
           incurs a retrieval cost.}
  \label{fig:adaptive_tradeoff}
\end{figure}%
}{}

We evaluate the adaptive adversary described in Section~\ref{sec:adaptive}
against \sad{} for all three attack types.  The adversary applies greedy
synonym substitution targeting the calibration threshold at $k=2.0$.

\paragraph{Key finding.}
Across all attacks, the adaptive adversary faces a fundamental tradeoff:
evasion gains (passages falling below the detection threshold) come at
the cost of retrieval degradation (ASR-R reduction).  Specifically:
\begin{itemize}
  \item \agentpoison{}: high evasion rate ($\approx 60\%$) but largest
    retrieval degradation ($\Delta\text{ASR-R} \approx 0.40$) because the
    centroid passage has extremely high initial similarity that requires
    many substitutions to reduce.
  \item \minja{}: moderate evasion rate ($\approx 40\%$) and moderate
    degradation ($\Delta\text{ASR-R} \approx 0.25$); bridging-step passages
    have intermediate initial similarity.
  \item \injecmem{}: lowest evasion rate ($\approx 20\%$) but also lowest
    initial anomaly signal; broad-anchor passages are already near the
    detection boundary.
\end{itemize}
These results confirm that \sad{} maintains meaningful protection even under
a white-box adaptive adversary: any successful evasion incurs a substantial
retrieval cost, reducing the threat to a level comparable to or lower than
the defended ASR-R of simpler defenses.

\subsection{Watermark Evasion}

We evaluate three evasion strategies against the unigram watermark
(Section~\ref{sec:defenses}): (i) paraphrasing at intensity 0.1--0.5,
(ii) copy-paste dilution with $r=0.5$--$3.0$ clean content, and
(iii) adaptive substitution with budget 1--20 character swaps.
Detailed results including z-score distribution curves are reported in
Appendix~\ref{sec:evasion_appendix}.

\IfFileExists{figures/p13_evasion_analysis.pdf}{%
\begin{figure}[t]
  \centering
  \includegraphics[width=0.82\linewidth]{p13_evasion_analysis.pdf}
  \caption{Watermark evasion analysis across three strategies:
           paraphrasing (varying synonym substitution intensity $\rho$),
           copy-paste dilution (dilution ratio $r$), and adaptive
           character substitution (budget $b$).
           $\tpr$ curves show that the unigram z-score defense
           remains robust ($\tpr > 0.80$) at low intensity/budget
           but degrades at aggressive evasion.}
  \label{fig:evasion_analysis}
\end{figure}%
}{}

\subsection{Ablation Studies}
\label{sec:ablation}

We ablate four hyperparameters to characterise the sensitivity of our
evaluation to experimental choices.  Full ablation tables are in
Appendix~\ref{sec:ablation_appendix}.

\paragraph{Corpus size.}
ASR-R does not decay with corpus size (50--500 entries), confirming that
poison entries act as extreme similarity outliers regardless of how many
benign entries are present.  This validates that our 200-entry synthetic
corpus reflects realistic deployment scales.

\paragraph{Retrieval top-\texorpdfstring{$k$}{k}.}
ASR-R increases monotonically with $k$ (1--20): a larger retrieval window
gives the adversary more opportunities for poison to appear in results.
At $k=1$, \agentpoison{} still achieves $\asrr > 0.9$ due to the trigger,
while \injecmem{} drops to $\approx 0.3$.

\paragraph{Poison count.}
\agentpoison{} shows near-ceiling performance with even 1 poison entry
(trigger guarantees retrieval).  \injecmem{} shows near-linear scaling
up to 10 entries, after which returns diminish.

\paragraph{\sad{} threshold \texorpdfstring{$k$}{k}.}
Lower $k$ yields higher TPR but also higher FPR (aggressive detection).
The default $k=2.0$ achieves TPR $\approx 0.85$ and FPR $\approx 0.05$
for \agentpoison{} (see Appendix~\ref{sec:ablation_appendix} and
Figure~\ref{fig:sad_ablation}).

\IfFileExists{figures/p13_ablation_sad_sigma.pdf}{%
\begin{figure}[t]
  \centering
  \begin{subfigure}[b]{0.48\linewidth}
    \includegraphics[width=\linewidth]{p13_ablation_sad_sigma.pdf}
    \caption{\sad{} threshold $k$ vs.\ TPR/FPR}
    \label{fig:sad_ablation}
  \end{subfigure}\hfill
  \begin{subfigure}[b]{0.48\linewidth}
    \IfFileExists{figures/p13_ablation_corpus.pdf}{%
      \includegraphics[width=\linewidth]{p13_ablation_corpus.pdf}%
    }{}
    \caption{Corpus size $N$ vs.\ $\asrr$}
    \label{fig:corpus_ablation}
  \end{subfigure}
  \caption{Ablation curves.
           Left: \sad{} threshold $k$ controls the TPR/FPR tradeoff;
           $k=2.0$ (vertical dashed) achieves $\tpr\!\approx\!0.85$,
           $\fpr\!\approx\!0.05$.
           Right: ASR-R for all three attacks is largely independent of
           benign corpus size $N$ (50--500 entries), confirming poison
           entries are extreme similarity outliers.}
  \label{fig:ablation_combined}
\end{figure}%
}{}
